% Generated draft: requires mathematical and source review.
\documentclass[11pt]{article}
\usepackage[margin=25mm]{geometry}
\usepackage{fontspec}
\setmainfont{Latin Modern Roman}
\usepackage{amsmath,amssymb,mathtools,unicode-math}
\setmathfont{Latin Modern Math}
\usepackage{enumitem,needspace,xurl,hyperref}
\setlength{\parindent}{0pt}
\setlength{\parskip}{5pt}
\setlength{\emergencystretch}{3em}
\newcommand{\N}{\mathbb{N}}
\newcommand{\Z}{\mathbb{Z}}
\newcommand{\Q}{\mathbb{Q}}
\newcommand{\R}{\mathbb{R}}
\newcommand{\C}{\mathbb{C}}
\newcommand{\F}{\mathbb{F}}
\newcommand{\A}{\mathbb{A}}
\newcommand{\PP}{\mathbb{P}}
\newcommand{\E}{\mathbb{E}}
\newcommand{\eps}{\varepsilon}
\newcommand{\dd}{\,\mathrm{d}}
\begin{document}
\section*{UnsolvedMath ID: 30002580}
\textbf{Source problem code:} catalogue entry
\subsection{Definitions and mathematical statement}
Let \(D\subset\C^n\) be a Cartan domain of rank \(r\), and let \(G=\operatorname{Aut}(D)\). The supplied record describes \(\mathcal R\) as the continuous-parameter algebra obtained from the \(K\)-radial parts of the \(G\)-invariant differential operators by allowing the dimension parameters to vary. It does not specify \(K\), the radial coordinates, the dimension parameters, their admissible range, or the construction of the resulting algebra. Thus this description identifies the intended object but does not determine it precisely.

The source poses the following questions about this generalized radial-part algebra:
\begin{enumerate}
\item Does \(\mathcal R\) admit a natural geometric realization?
\item Is \(\mathcal R\) commutative, meaning that \(AB=BA\) for every \(A,B\in\mathcal R\)?
\item Are all members of \(\mathcal R\) normal operators with respect to the generalized radial measure?
\item Does \(\mathcal R\) admit sufficiently many joint eigenfunctions to define an analogue of the Harish--Chandra transform?
\item Is every member of \(\mathcal R\) a canonical function of a distinguished generating set?
\end{enumerate}

For clarity, once a radial space \(X\), a positive measure \(\mu\), and operator realizations on \(H=L^2(X,\mu)\) have been specified, normality of an operator \(A\) means that \(A\) is closed and densely defined and that \(A^*A=AA^*\), including equality of the domains of these products. The supplied record does not specify these data or say whether it instead intends formal normality of differential expressions. Likewise, algebraic commutativity of radial differential expressions must be distinguished from commutation properties of their Hilbert-space realizations.

A joint eigenfunction would satisfy \(A\phi=\lambda_A\phi\) for every \(A\in\mathcal R\), in a function space and with an interpretation of these equations that remain unspecified. The required completeness or spectral expansion, the proposed transform, the meaning of a natural geometric realization, and the functional calculus intended by canonical functions of generators are also unspecified. Consequently, the five questions above preserve the source's intended targets, but a fully quantified conjecture requires additional definitions and hypotheses; no particular parameter range, boundary condition, operator extension, or functional calculus is selected here.
\subsection{Short English statement}
For a Cartan domain, the problem asks whether extending radial parts of invariant differential operators to continuous dimension parameters yields an algebra with a geometric interpretation, commutativity, normal operator realizations, enough joint eigenfunctions for a Harish--Chandra type transform, and canonical generators with an associated functional calculus. The supplied record leaves the parameter range and essential analytic definitions unresolved, so this is a qualified statement-only draft.
\subsection{Sources}
\begin{itemize}
\item Local UnsolvedMath registry: \texttt{data/unsolvedmath/problems.json}, permanent ID \(30002580\), source problem code \texttt{catalogue entry}, title ``Geometry and Spectral Theory of Generalized Radial-Part Algebras.'' The supplied record belongs to the Oberwolfach Reports Open Problems collection and gives the proposed year as 2014; no individual proposer is named.
\item Supplied source citation: \emph{Hilbert Modules and Complex Geometry} (2015), Oberwolfach Reports, DOI 10.4171/owr/2014/21, pp. 1139--1219. \url{https://doi.org/10.4171/owr/2014/21}. The background describes these reports as documentation of mathematical workshops and their open-problem discussions. The supplied citation's 2015 date, its DOI containing 2014, and the proposed year 2014 are retained without resolving their relationship.
\item Supplied literature source, labeled ``One-parameter symmetric-domain calculus'': \url{https://doi.org/10.1002/mana.200610527}. No authors, full bibliographic title, or precise supporting result are supplied for this reference.
\item Literature assessment dated exactly \texttt{2026-08-21}: the supplied curation labels the problem open and reports spectral and functional-calculus results for rank-one generalized radial operators, while describing the general geometric realization, commutativity, normality, joint eigenfunctions, and canonical generators for the continuous-parameter family as unresolved. Both references above are listed among the sources checked by that curation. This is supplied dated curation, not an independently verified literature assessment; the background expressly requests independent verification of time-sensitive status and cited claims.
\item The supplied registry labels the statement \texttt{corrected\_verified} and says that mathematical content was preserved while notation, typography, or omitted definitions were repaired for publication. This is a supplied editorial qualification, not independent verification of the construction or mathematical claims in this draft.
\end{itemize}
\end{document}
